\subsection{\texorpdfstring{\(G_{2}\) 型系统}{G\_2 型系统}}

\begin{enumerate}
\def\labelenumi{(\Roman{enumi})}
\tightlist
\item
  令 E 为 \(R^{3}\) 中方程为
\end{enumerate}

\[
\xi_ {1} + \xi_ {2} + \xi_ {3} = 0.
\]

  令 \(R\) 为由 \(\alpha \in L_0 \cap V\) 且满足 \((\alpha | \alpha) = 2\) 或 \((\alpha | \alpha) = 6\) 的元素组成的集合。\(R\) 的元素为

\[
\begin{array}{c} \pm (\varepsilon_ {1} - \varepsilon_ {2}), \quad \pm (\varepsilon_ {1} - \varepsilon_ {3}), \quad \pm (\varepsilon_ {2} - \varepsilon_ {3}), \quad \pm (2 \varepsilon_ {1} - \varepsilon_ {2} - \varepsilon_ {3}), \\ \pm (2 \varepsilon_ {2} - \varepsilon_ {1} - \varepsilon_ {3}), \quad \pm (2 \varepsilon_ {3} - \varepsilon_ {1} - \varepsilon_ {2}). \end{array}
\]

  于是，R 生成 V，且 \(\frac{2(\alpha|\beta)}{(\beta|\beta)}\in\mathbf{Z}\) 对所有 \(\alpha,\beta\in R\) 成立：当 \(\beta=\pm(\varepsilon_{i}-\varepsilon_{j})\) 且 \(i\neq j\) 时这显然成立；例如，当 \(\beta=2\varepsilon_{1}-\varepsilon_{2}-\varepsilon_{3}\) 时，我们有 \((\alpha|\beta)\in3\mathbf{Z}\) ，断言同样成立。因此，R 是 V 中的约化根系。根的数目为 n=12。

\begin{enumerate}
\def\labelenumi{(\Roman{enumi})}
\setcounter{enumi}{1}
\tightlist
\item
  置 \(\alpha_{1} = \varepsilon_{1} - \varepsilon_{2},\alpha_{2} = -2\varepsilon_{1} + \varepsilon_{2} + \varepsilon_{3}\) 。则这些根为
\end{enumerate}

\[
\begin{array}{r l} & {\pm \alpha_ {1}, \quad \pm (\alpha_ {1} + \alpha_ {2}), \quad \pm (2 \alpha_ {1} + \alpha_ {2}), \quad \pm \alpha_ {2},} \\ & {\qquad \pm (3 \alpha_ {1} + \alpha_ {2}), \quad \pm (3 \alpha_ {1} + 2 \alpha_ {2}).} \end{array}
\]

因此，\((\alpha_{1},\alpha_{2})\) 是 \(R\) 的一个基。我们有 \(\| \alpha_{1}\|^{2} = 2,\| \alpha_{2}\|^{2} = 6.\) \((\alpha_{1}|\alpha_{2}) = -3,\) 所以 \(R\) 是 \(G_{2}\) 型系统。正根为 \(\alpha_{1},\alpha_{1} + \alpha_{2},2\alpha_{1} + \alpha_{2},\) \(3\alpha_{1} + \alpha_{2},3\alpha_{1} + 2\alpha_{2}\) 。

\begin{enumerate}
\def\labelenumi{(\Roman{enumi})}
\setcounter{enumi}{2}
\item
  我们有 \(h = \frac{n}{2} = 6\) 。
\item
  最高根为 \(\tilde{\alpha} = 3\alpha_{1} + 2\alpha_{2} = -\varepsilon_{1} - \varepsilon_{2} + 2\varepsilon_{3}\) 。我们有 \((\tilde{\alpha}|\alpha_{1}) = 0, (\tilde{\alpha}|\alpha_{2}) = 3\) 。扩充的 Dynkin 图为
\end{enumerate}

\[
\alpha_ {1} \quad \alpha_ {2}
\]

\begin{enumerate}
\def\labelenumi{(\Alph{enumi})}
\setcounter{enumi}{21}
\tightlist
\item
  逆系统是下列向量集：
\end{enumerate}

\[
\begin{array}{l} \pm \alpha_ {1}, \quad \pm (\alpha_ {1} + \alpha_ {2}), \quad \pm (2 \alpha_ {1} + \alpha_ {2}), \quad \pm \frac {1}{3} \alpha_ {2}, \\ \qquad \pm \frac {1}{3} (3 \alpha_ {1} + \alpha_ {2}), \quad \pm \frac {1}{3} (3 \alpha_ {1} + 2 \alpha_ {2}). \end{array}
\]

  有 10 个根不与 \(\alpha_{1}\) 正交；其中 4 个根满足 \(n(\beta, \alpha_{1}) = \pm 1\) ，另外 4 个满足 \(n(\beta, \alpha_{1}) = \pm 3\) ，另有 \(n(\beta, \alpha_{1}) = \pm 2\) 的根满足 \(\beta = \pm \alpha_{1}\) 。因此，\(\alpha_{1}\) 在 \(\Phi_{\mathbb{R}}\) 中的长度平方为 \(4(4.1 + 4.9 + 2.4)^{-1} = \frac{1}{12}\) 。从而 \(\Phi_{\mathbb{R}}(x, y) = (x|y)/24\) 。现在对 \(x = y = \alpha_{1}\) 应用 §1 第 12 号的公式 (18)，得到

\[
2 + 4. \frac {1}{4} + 4. \frac {1}{4} = \gamma (\mathrm{R}). \frac {1}{12}
\]

所以 \(\gamma (\mathbb{R}) = 48\)

\begin{enumerate}
\def\labelenumi{(\Roman{enumi})}
\setcounter{enumi}{5}
\tightlist
\item
  (VI) 和 (VII) 正根之和的一半为
\end{enumerate}

\[
\rho = 5 \alpha_ {1} + 3 \alpha_ {2}.
\]

基本权 \(\omega_{1}\) 和 \(\omega_{2}\) 分别与 \(\alpha_{2}\) 和 \(\alpha_{1}\) 正交，因而分别与 \(2\alpha_{1} + \alpha_{2}\) 和 \(3\alpha_{1} + 2\alpha_{2}\) 成比例。我们有

\[
\omega_ {1} + \omega_ {2} = \rho = 5 \alpha_ {1} + 3 \alpha_ {2} = (2 \alpha_ {1} + \alpha_ {2}) + (3 \alpha_ {1} + 2 \alpha_ {2}).
\]

因此，

\[
\omega_ {1} = 2 \alpha_ {1} + \alpha_ {2}, \quad \omega_ {2} = 3 \alpha_ {1} + 2 \alpha_ {2} = \tilde {\alpha}.
\]

\begin{enumerate}
\def\labelenumi{(\Roman{enumi})}
\setcounter{enumi}{7}
\item
  例如，Q(R) 由 \(\varepsilon_{1}-\varepsilon_{2}\) 和 \(\varepsilon_{1}-\varepsilon_{3}\) 生成。由 (VI) 和 (VII)，P(R) = Q(R)。连通指标为 1。
\item
  指数族有 2 项：由于 1 和 \(h - 1 = 5\) 是指数，它们就是全部指数。
\item
  我们有 \((\alpha_{1},\alpha_{2}) = \frac{5\pi}{6}\) ，所以 \(\mathrm{W}(\mathbb{R})\) 同构于阶为 12 的二面体群。
\item
  Dynkin 图的唯一自同构是恒等自同构，所以 \(A(R) = W(R)\) 且 \(w_{0} = -1\) 。
\end{enumerate}

